\section{Related Work}

Multi-objective optimization has long studied how to compare and select solutions when several objectives must be considered simultaneously. Classical formulations describe solution quality through objective vectors and use Pareto dominance to identify solutions that represent meaningful trade-offs among conflicting objectives \cite{miettinen1999nonlinear}. Evolutionary multi-objective algorithms, such as NSGA-II, further popularized Pareto-front-based reasoning by maintaining diverse sets of nondominated solutions during optimization \cite{deb2002fast}. These foundations are important for this project because Gen-MOR also relies on the distinction between nondominated and dominated plans. However, rather than directly solving a multi-objective optimization problem, Gen-MOR studies whether neural generative models can learn and generate objective-vector plans while preserving their dominance-based structure.

Generative modeling provides the learning foundation for this work. Autoencoders learn compressed representations by mapping inputs into a latent space and reconstructing them from that representation \cite{hinton2006reducing}. Variational autoencoders extend this idea by learning a probabilistic latent space, making them suitable for sampling and generation \cite{kingma2014autoencoding}. Conditional generative models further allow generation to be guided by auxiliary information, such as class labels or desired attributes \cite{sohn2015learning}. More broadly, the generative modeling literature includes adversarial models, normalizing flows, and diffusion models, each of which provides a different way to learn and sample from complex data distributions \cite{goodfellow2014generative,papamakarios2021normalizing,ho2020denoising}. Gen-MOR uses this general family of ideas, but applies them to a structured multi-objective planning setting where validity depends on dominance relationships rather than only on visual, chemical, or statistical realism.

Recent work has increasingly connected generative modeling with multi-objective optimization. Abeer et al. proposed a multi-objective latent space optimization method for generative molecular design models \cite{abeer2024multiobjective}. Their approach uses Pareto-efficiency-based weighting during iterative retraining so that the generative model becomes more likely to produce molecules with improved multi-objective trade-offs. This work is relevant because it demonstrates that Pareto information can be used to influence a learned latent generative model. However, its goal is to improve molecular candidates with respect to multiple properties, while Gen-MOR focuses on whether generated plan vectors preserve dominance labels and structural consistency.

Latent Pareto set learning provides another closely related direction. The MLPS framework studies multi-objective molecular design by learning a mapping from preference vectors to Pareto-optimal solutions in a continuous latent space \cite{mlps2024latentparetoset}. This approach addresses the limitation of producing only a small number of Pareto candidates by attempting to learn a broader representation of the Pareto set. Gen-MOR is related to this idea because it also treats latent representation learning as useful for multi-objective reasoning. The difference is that Gen-MOR does not primarily aim to approximate the entire Pareto set. Instead, it evaluates whether generated samples remain consistent with the nondominated and dominated plan classes produced by a MORNP-style setting.

More recent generative approaches have used flow-based and diffusion-based models to guide samples toward Pareto-optimal regions. ParetoFlow introduces guided flow matching for offline multi-objective optimization \cite{yuan2025paretoflow}. It uses multi-objective predictor guidance and weight vectors to encourage generated samples to cover different regions of the Pareto front. Similarly, PROUD formulates multi-objective generation as a constrained diffusion process, where generated samples should maintain generation quality while moving toward Pareto optimality \cite{yao2024proud}. These methods show that advanced generative models can be explicitly guided by multi-objective structure. Gen-MOR differs from them in evaluation emphasis: instead of only asking whether generated samples approach the Pareto front, Gen-MOR asks whether generated samples preserve their intended dominance category as $P^{+}$ or $P^{-}$.

Pareto-aware generation has also been studied in scientific design problems such as proteins and molecules. MosPro proposes a machine-learning-guided framework for multi-objective protein sequence design, using Pareto-optimal sampling to generate sequences that balance multiple protein properties \cite{luo2025mospro}. This line of work shows that Pareto-aware generation is practically important in domains where a design must satisfy several competing scientific objectives. Antoniuk et al. further show that static generative models may fail to advance the known Pareto front unless generation is combined with active learning, external verification, and iterative retraining \cite{antoniuk2026active}. This observation is especially relevant to Gen-MOR because it reinforces the need to verify generated samples using a domain-specific structural criterion rather than relying only on ordinary learning metrics.

Taken together, existing work shows that generative models can support multi-objective search, latent-space optimization, Pareto-front approximation, and scientific design. However, most of these methods focus on generating samples that improve objective values or approximate the Pareto front. Gen-MOR addresses a complementary problem: generating multi-objective plan vectors that preserve dominance-based labels. In this project, a generated $P^{+}$ sample should remain nondominated, while a generated $P^{-}$ sample should remain dominated under the selected objective subset. This motivates the use of Dominance Validity Rate (DVR) as a structural evaluation metric alongside reconstruction error, classification accuracy, and F1 score. Therefore, Gen-MOR positions itself not as another Pareto-front optimizer, but as a dominance-preserving generative modeling framework for MORNP-style multi-objective planning data.